%%%%%%%%%%%%%%%%%%%%%%%%%%%%%%%%%%%%%%%%%%%%%%%%%%%%%%%%%%%%%%%%%%%%%%%%
% Plantilla TFG/TFM
% Universidad de Murcia. Facultad de Informática
% Realizado por: José Manuel Requena Plens
% Modificado: Pablo José Rocamora Zamora
% Contacto: pablojoserocamora@gmail.com
%%%%%%%%%%%%%%%%%%%%%%%%%%%%%%%%%%%%%%%%%%%%%%%%%%%%%%%%%%%%%%%%%%%%%%%%


\chapter*{Resumen}
\addcontentsline{toc}{chapter}{Resumen} %

La crio-microscopía electrónica de tomografía (\textit{cryo-ET}) es una técnica de imagen de alta
resolución que permite estudiar la organización tridimensional de macromoléculas en su entorno
celular nativo. Mediante la adquisición de proyecciones 2D desde distintos ángulos de inclinación y su posterior reconstrucción computacional, se obtienen tomogramas 3D que revelan la arquitectura interna de células y tejidos a escala nanométrica. El análisis automatizado de estos tomogramas mediante redes neuronales profundas requiere grandes volúmenes de datos de entrenamiento etiquetados, cuya obtención experimental es costosa y laboriosa. Una alternativa consolidada es la generación de \textbf{tomogramas sintéticos}: volúmenes simulados con distribuciones realistas de proteínas, membranas y ruido, que reproducen las características estadísticas de los datos reales.\\

En este contexto se enmarca el \textit{framework} \textbf{PolNet}~\cite{martinezSanchezPolnetGithub}, un marco de simulación de tomogramas sintéticos que incorpora el módulo \textit{SAWLC} (\textit{Self-Avoiding Worm-Like Chain}) para la
inserción de proteínas en el volumen de interés. Aunque SAWLC produce distribuciones
biológicamente plausibles, su rendimiento computacional limita su aplicabilidad cuando se
requieren grandes colecciones de datos de entrenamiento con un grado de empaquetamiento alto o se trabaja con proteínas de gran variedad.\\

El presente Trabajo Fin de Grado aborda el diseño, implementación y evaluación de un algoritmo
alternativo de inserción de proteínas denominado \textbf{Tetris}~\cite{Hamo2025tetris_rep} basado en el \textit{paper Template Learning}~\cite{harastani2025template}, desarrollado en un repositorio
independiente con el objetivo de compararlo frente a SAWLC y valorar su viabilidad como reemplazo.
Tetris emplea correlación cruzada 3D basada en la transformada de Fourier rápida (FFT) para
identificar, en cada paso, la posición óptima de inserción de una proteína respecto a las ya
colocadas, promoviendo un empaquetamiento denso y espacialmente coherente. Además, se desarrolla una versión acelerada por GPU.\\

Los principales objetivos del trabajo son: (i)~implementar el algoritmo Tetris,
(ii)~desarrollar su versión GPU, (iii)~diseñar un banco de pruebas que evalúe el rendimiento frente a SAWLC en términos de ocupancia volumétrica, número de monómeros insertados y tiempo de ejecución, y (iv)~analizar el impacto de la aceleración GPU sobre dichas métricas.\\ 
 

Los resultados obtenidos muestran que Tetris logra una ocupancia superior a la de SAWLC en igualdad de condiciones, lo que lo posiciona como una alternativa competitiva para la generación eficiente de datos sintéticos en cryo-ET.
